\chapter{Perimeter and Area}\label{ch:g5:areas}

Grade 4 counted grid squares (\cref{ch:g4:measure}); this chapter
earns the first \omterm{def:g4:measure:area}{area} \emph{formula} --- length times width for the
rectangle --- and learns the units cm$^2$ and m$^2$. The formulas for
triangles and other shapes ripen in \cref{ch:g6:measure} and
\cref{ch:g7:areas}.

\section{Two different measures}

\begin{example}[Border vs surface]\label{ex:g5:areas:contrast}
A gardener needs \emph{\omterm{def:g4:measure:perimeter}{perimeter}} to buy the fence and \emph{\omterm{def:g4:measure:area}{area}}
to buy the grass seed. The two do not follow each other: stretching
a rectangle thinner and longer can keep its \omterm{def:g4:measure:area}{area} while its \omterm{def:g4:measure:perimeter}{perimeter}
grows --- compare a $6 \times 2$ and a $4 \times 3$ rectangle
(same \omterm{def:g4:measure:area}{area} $12$, \omterm{def:g4:measure:perimeter}{perimeters} $16$ and $14$).
\end{example}

\section{The rectangle formula}

\begin{proposition}[Area of a rectangle]\label{prop:g5:areas:rectangle}
A rectangle of length $L$ and width $w$ (in the same unit) has \omterm{def:g4:measure:area}{area}
\[
A = L \times w ,
\]
in square units: cm$^2$ (squares of side $1$ cm), m$^2$ (squares of
side $1$ m), \dots
\end{proposition}

\begin{proof}[Why, by counting]
Cover the rectangle with unit squares: $w$ rows of $L$ squares each,
so $L \times w$ squares in total --- exactly the multiplication
rectangle of \cref{ch:g3:mult}.
\end{proof}

\begin{omfigure}
\begin{tikzpicture}[scale=0.62]
  \draw[gray!50, thin] (0,0) grid (6,3);
  \draw[very thick, omDef] (0,0) rectangle (6,3);
  \fill[omThm!30] (0,2) rectangle (1,3);
  \draw[omThm, thick] (0,2) rectangle (1,3);
  \node[right, font=\small, omThm] at (1.15,2.5) {$1$ cm$^2$};
  \node[below, font=\small] at (3,-0.15) {$6$ cm};
  \node[left, font=\small] at (0,1.5) {$3$ cm};
  \node[font=\small] at (3.9,1.2) {$6 \times 3 = 18$ cm$^2$};
\end{tikzpicture}

{\small Three rows of six centimeter-squares: the formula
$L \times w$ is the counting, written once and for all.}
\end{omfigure}

\begin{example}\label{ex:g5:areas:rectangle}
A rug measures $2.5$ m by $2$ m: \omterm{def:g4:measure:area}{area} $2.5 \times 2 = 5$ m$^2$.
A stamp measures $3$ cm by $2.4$ cm: \omterm{def:g4:measure:area}{area} $3 \times 2.4 = 7.2$
cm$^2$. Same formula, any unit --- as long as both sides use the
\emph{same} one.
\end{example}

\begin{example}[Square units convert by 100]\label{ex:g5:areas:units}
$1$ m $= 100$ cm, but $1$ m$^2 = 100 \times 100 = 10\,000$ cm$^2$:
a square meter is a $100$-by-$100$ grid of square centimeters. \omterm{def:g4:measure:area}{Area}
units jump by hundreds, not tens.
\end{example}

\section{Composite figures}

\begin{method}[Cut, add, subtract]\label{met:g5:areas:composite}
For a figure made of rectangles (an L, a T, a frame):
\begin{enumerate}
  \item cut it into rectangles, or complete it into a big rectangle;
  \item compute each rectangular \omterm{def:g4:measure:area}{area} with the formula;
  \item add the pieces --- or subtract the hole from the big
        rectangle;
  \item check against a rough count of grid squares.
\end{enumerate}
\end{method}

\begin{example}\label{ex:g5:areas:composite}
An L-shaped room: a $6$ m $\times$ $4$ m rectangle with a
$2$ m $\times$ $2$ m corner missing.
\[
\text{Area} = 6 \times 4 - 2 \times 2 = 24 - 4 = 20 \text{ m}^2 .
\]
Or cut the L into a $6 \times 2$ strip and a $4 \times 2$ strip:
$12 + 8 = 20$ m$^2$ --- two roads, one answer.
\end{example}

\begin{example}[Half a rectangle]\label{ex:g5:areas:halfrect}
Cutting a rectangle along a diagonal gives two triangles of the same
\omterm{def:g4:measure:area}{area}: each is \emph{\omterm{def:g3:division:half}{half}} the rectangle. A right triangle with legs
$6$ cm and $4$ cm therefore has \omterm{def:g4:measure:area}{area} $\frac{6 \times 4}{2} = 12$
cm$^2$ --- a picture worth remembering for
\cref{ch:g6:measure}, where it becomes a formula.
\end{example}

\begin{omfigure}
\begin{tikzpicture}[scale=0.62]
  \draw[gray!50, thin] (0,0) grid (6,4);
  \draw[very thick] (0,0) rectangle (6,4);
  \fill[omDef!25] (0,0) -- (6,0) -- (0,4) -- cycle;
  \draw[very thick, omDef] (6,0) -- (0,4);
  \node[font=\small] at (1.7,1.2) {half};
  \node[font=\small] at (4.3,2.8) {half};
\end{tikzpicture}

{\small The diagonal cuts the $6 \times 4$ rectangle into two equal
triangles of $12$ squares each.}
\end{omfigure}

\section{Exercises}

\begin{exercise}[$\star$]\label{exo:g5:areas:1}
Compute the \omterm{def:g4:measure:area}{area} and the \omterm{def:g4:measure:perimeter}{perimeter} of a rectangle $8$ cm $\times$
$5$ cm. Which answer is in cm, which in cm$^2$?
\end{exercise}

\begin{exercise}[$\star$]\label{exo:g5:areas:2}
Compute the \omterm{def:g4:measure:area}{areas}: a square of side $9$ cm; a rectangle
$12$ m $\times$ $7$ m; a rectangle $4.5$ cm $\times$ $6$ cm.
\end{exercise}

\begin{exercise}[$\star$]\label{exo:g5:areas:3}
A rectangle has \omterm{def:g4:measure:area}{area} $63$ cm$^2$ and length $9$ cm. Find its width,
then its \omterm{def:g4:measure:perimeter}{perimeter}.
\end{exercise}

\begin{exercise}[$\star$]\label{exo:g5:areas:4}
Draw two different rectangles with \omterm{def:g4:measure:area}{area} $24$ squares on grid paper,
and compute both \omterm{def:g4:measure:perimeter}{perimeters}. Same \omterm{def:g4:measure:area}{area} --- same \omterm{def:g4:measure:perimeter}{perimeter}?
\end{exercise}

\begin{exercise}[$\star$]\label{exo:g5:areas:5}
Convert: $3$ m$^2$ in cm$^2$; \quad $50\,000$ cm$^2$ in m$^2$.
(Remember \cref{ex:g5:areas:units}: by hundreds!)
\end{exercise}

\begin{exercise}[$\star$]\label{exo:g5:areas:6}
A T-shaped figure is made of a $8 \times 2$ horizontal bar on top of
a $2 \times 5$ vertical bar (in cm). Compute its \omterm{def:g4:measure:area}{area} by adding two
rectangles.
\end{exercise}

\begin{exercise}[$\star$]\label{exo:g5:areas:7}
A picture frame: a $30$ cm $\times$ $20$ cm rectangle with a
$24$ cm $\times$ $14$ cm rectangular window cut out. What \omterm{def:g4:measure:area}{area} of
wood does the frame use?
\end{exercise}

\begin{exercise}[$\star$]\label{exo:g5:areas:8}
Compute the \omterm{def:g4:measure:area}{area} of a right triangle with legs $8$ cm and $6$ cm
(\omterm{def:g3:division:half}{half} a rectangle, \cref{ex:g5:areas:halfrect}).
\end{exercise}

\begin{exercise}[$\star$]\label{exo:g5:areas:9}
A rectangular vegetable patch measures $7$ m by $4$ m. Seed costs
$2$ per square meter, and fencing $3$ per meter. Compute the cost of
the seed, then the cost of the fence.
\end{exercise}

\begin{exercise}[$\star\star$]\label{exo:g5:areas:10}
A corridor floor is $12$ m long and $2$ m wide, and must be covered
with square tiles of side $50$ cm. How many tiles are needed?
(Convert first, or count tiles along each direction.)
\end{exercise}

\begin{exercise}[$\star\star$]\label{exo:g5:areas:11}
Double the sides of a $4$ cm $\times$ $3$ cm rectangle. What happens
to its \omterm{def:g4:measure:perimeter}{perimeter}? To its \omterm{def:g4:measure:area}{area}? (Compute both before and after ---
the two answers differ!)
\end{exercise}
